\documentclass[10pt]{extarticle}
\usepackage[a4paper,margin=1.5cm]{geometry}
\usepackage[T1]{fontenc}
\usepackage{amsmath,amssymb}
\usepackage{graphicx}
\usepackage[export]{adjustbox}
\usepackage{multicol}
\usepackage{enumitem}
\usepackage{booktabs,array,tabularx}
\setlength{\columnsep}{0.7cm}
\setlength{\parindent}{0pt}
\setlength{\parskip}{3pt}
\pagestyle{empty}
\makeatletter
% figure registry (figdefs.tex is generated by the extract step)
\newcommand{\qfigdef}[4]{\@namedef{qf@file@#1}{#2}\@namedef{qf@w@#1}{#3}\@namedef{qf@wide@#1}{#4}}
\newcommand{\qf@check}[1]{\@ifundefined{qf@file@#1}{\PackageError{score2}{Unknown figure id #1}{}}{}}
\newcommand{\qf@img}[4][]{\includegraphics[width=#2\dimexpr\@nameuse{qf@w@#4}pt\relax,max width=#3,max height=0.45\textheight,#1]{\@nameuse{qf@file@#4}}}
% Full-width questions: finalize wraps every question that contains a wide figure in \qfullbegin...\qfullend,
% so the WHOLE question (stem, figures, options) is set across the page in reading order, then the two columns
% resume. The question's head - its number, stem and first figure (\qfullsplit marks the end of the head) -
% is kept on one page, so a number is never left at a page foot with its figure overleaf; the rest may break
% across pages normally, so no big gaps are left. Without a \qfullsplit, the whole question is kept together
% if it is under 45% of a page. Never write these yourself.
\newif\ifqf@full
\newif\ifqf@boxed
\newsavebox{\qf@box}
\newcommand{\qf@open}{\global\qf@fulltrue\par\global\qf@boxedtrue
  \setbox\qf@box\vbox\bgroup\hsize\textwidth\linewidth\textwidth}
\newcommand{\qf@emit}[1]{\par\egroup\global\qf@boxedfalse
  \ifdim\dimexpr\ht\qf@box+\dp\qf@box\relax<#1\textheight\noindent\usebox{\qf@box}\par\else\unvbox\qf@box\fi}
\newcommand{\qfullsplit}{\ifqf@boxed\qf@emit{0.85}\fi}
\newcommand{\qf@close}{\ifqf@boxed\qf@emit{0.45}\fi\global\qf@fullfalse}
\newcommand{\qfullbegin}{\end{multicols}\qf@open}
\newcommand{\qfullend}{\qf@close\begin{multicols}{2}}
\newcommand{\qfullstart}{\qf@open}
\newcommand{\qfulldone}{\qf@close}
% \qfig[scale]{id}: block figure, centred, directly where it is written (size set by the tool).
% Fallback if a wide figure is not inside a full-width question: break out of the columns in place - no floats.
\newcommand{\qfig}[2][1]{\qf@check{#2}%
  \ifqf@full\qf@blk{#1}{#2}\else
  \ifnum\@nameuse{qf@wide@#2}=1 \expandafter\@firstoftwo\else\expandafter\@secondoftwo\fi
  {\end{multicols}\begin{center}\qf@img{#1}{\textwidth}{#2}\end{center}\begin{multicols}{2}}%
  {\qf@blk{#1}{#2}}\fi}
\newcommand{\qf@blk}[2]{\par\smallskip{\centering\qf@img{#1}{\linewidth}{#2}\par}\smallskip}
% \qfigi[scale]{id}: inline figure (inside option lists / tables / minipages); never breaks columns.
\newcommand{\qfigi}[2][1]{\qf@check{#2}\qf@img[valign=c]{#1}{\linewidth}{#2}}
\makeatother
% structure macros - lint reads these to check completeness and figure placement
\newcommand{\papersec}[1]{\par\bigskip{\centering\Large\bfseries PAPER-#1\par}\medskip}
\newcommand{\partsec}[2]{\par\medskip\noindent{\large\bfseries PART-#1: #2}\par\smallskip}
\newcommand{\sect}[1]{\par\smallskip\noindent\textbf{#1}\par}
\newcommand{\qn}[1]{\par\smallskip\noindent\textbf{#1.}~}
\newcommand{\paraq}[2]{\par\smallskip\noindent\textbf{Paragraph for Questions #1 and #2}\par}
\newenvironment{qoptions}{\begin{itemize}[leftmargin=1.8em,itemsep=1pt,topsep=2pt,label={}]}{\end{itemize}}
\newcommand{\opt}[2]{\item[(#1)] #2}
% figure registry (generated by score2_tools.py finalize)
\qfigdef{p1_phy_s0_q01_1}{SRGBT - Test 13 (17 Oct 2021)-img/p1_phy_s0_q01_1.png}{66.8}{0}
\qfigdef{p1_phy_s0_q01_A}{SRGBT - Test 13 (17 Oct 2021)-img/p1_phy_s0_q01_A.png}{58.7}{0}
\qfigdef{p1_phy_s0_q01_3}{SRGBT - Test 13 (17 Oct 2021)-img/p1_phy_s0_q01_3.png}{58.6}{0}
\qfigdef{p1_phy_s0_q01_4}{SRGBT - Test 13 (17 Oct 2021)-img/p1_phy_s0_q01_4.png}{58.2}{0}
\qfigdef{p1_phy_s0_q01_5}{SRGBT - Test 13 (17 Oct 2021)-img/p1_phy_s0_q01_5.png}{58.7}{0}
\qfigdef{p1_phy_s0_q02_1}{SRGBT - Test 13 (17 Oct 2021)-img/p1_phy_s0_q02_1.png}{45.9}{0}
\qfigdef{p1_phy_s0_q03_1}{SRGBT - Test 13 (17 Oct 2021)-img/p1_phy_s0_q03_1.png}{154.7}{0}
\qfigdef{p1_phy_s0_q05_1}{SRGBT - Test 13 (17 Oct 2021)-img/p1_phy_s0_q05_1.png}{119.4}{0}
\qfigdef{p1_phy_s0_q06_1}{SRGBT - Test 13 (17 Oct 2021)-img/p1_phy_s0_q06_1.png}{105.6}{0}
\qfigdef{p1_phy_s0_q08_1}{SRGBT - Test 13 (17 Oct 2021)-img/p1_phy_s0_q08_1.png}{116.2}{0}
\qfigdef{p1_phy_s0_q10_1}{SRGBT - Test 13 (17 Oct 2021)-img/p1_phy_s0_q10_1.png}{88.3}{0}
\qfigdef{p1_phy_s0_q11_1}{SRGBT - Test 13 (17 Oct 2021)-img/p1_phy_s0_q11_1.png}{164.4}{0}
\qfigdef{p1_phy_s0_q12_1}{SRGBT - Test 13 (17 Oct 2021)-img/p1_phy_s0_q12_1.png}{101.2}{0}
\qfigdef{p1_phy_s0_q13_1}{SRGBT - Test 13 (17 Oct 2021)-img/p1_phy_s0_q13_1.png}{81.0}{0}
\qfigdef{p1_phy_s0_q14_2}{SRGBT - Test 13 (17 Oct 2021)-img/p1_phy_s0_q14_2.png}{115.3}{0}
\qfigdef{p1_phy_s0_q16_1}{SRGBT - Test 13 (17 Oct 2021)-img/p1_phy_s0_q16_1.png}{96.1}{0}
\qfigdef{p1_phy_s0_q18_1}{SRGBT - Test 13 (17 Oct 2021)-img/p1_phy_s0_q18_1.png}{130.9}{0}
\raggedcolumns
\makeatletter\@beginparpenalty=10000\makeatother
\newcolumntype{L}[1]{>{\raggedright\arraybackslash}p{#1}}
\begin{document}
\begin{center}{\Large\bfseries SRGBT --- Test 13 (17 Oct 2021)}\end{center}
\begin{multicols}{2}
\raggedcolumns
\papersec{1}
\partsec{1}{PHYSICS}
\qn{1} The piston is massless and the spring is ideal and initially stretched the piston cylinder arrangement encloses an ideal gas. If the gas is heated quasistatically the PV graph\par
\qfig{p1_phy_s0_q01_1}
\begin{qoptions}
\opt{A}{\qfigi{p1_phy_s0_q01_A}}
\opt{B}{\qfigi{p1_phy_s0_q01_3}}
\opt{C}{\qfigi{p1_phy_s0_q01_4}}
\opt{D}{\qfigi{p1_phy_s0_q01_5}}
\end{qoptions}
\qn{2} There is a sample of air at standard conditions in the vertical cylinder shown in the figure. The piston moves without friction and its mass is same as that of the pan on the other side of the pulley. Some fixed amount of sand is slowly poured either to the top of the piston or into the pan. The temperature remains the same. Then:\par
\qfig{p1_phy_s0_q02_1}
\begin{qoptions}
\opt{A}{Displacement of the piston will be greater if sand is poured in the pan}
\opt{B}{Displacement of the piston will be greater \ if sand is poured on to the top of the piston.}
\opt{C}{Displacement of the piston will be same in both cases}
\opt{D}{Displacement of piston is greater or lesser in any case, will depend on the initial volume \& pressure of the gas}
\end{qoptions}
\qn{3} A monoatomic ideal gas undergoes a process consists of concentric semi-circles on PV graph, starting from A(P$_0$, V$_0$) and ending at B(P$_0$, 2V$_0$). The point B is the common centre of the \ PV semi circles. Let r$_\mathrm n$ be the magnitude of radius of nth circle. It is given that r$_\mathrm n$ = 2/3 r$_{\mathrm n-1}$, n $\in$ N. Find the magnitude of heat gained (or rejected) by the system in undergoing whole process from A to B. Process is isobaric from (n $-$ 1)$^\mathrm{th}$ to n$^\mathrm{th}$ circle.\par
\qfig{p1_phy_s0_q03_1}
\begin{qoptions}
\opt{A}{$\left(\dfrac{65-9\pi}{26}\right)\mathrm P_0\mathrm V_0$}
\opt{B}{$\left(\dfrac{4+\pi}2\right)\mathrm P_0\mathrm V_0$}
\opt{C}{$\left(\dfrac{7+\pi}2\right)\mathrm P_0\mathrm V_0$}
\opt{D}{$\left(\dfrac{65+9\pi}{26}\right)\mathrm P_0\mathrm V_0$}
\end{qoptions}
\qn{4} A solid object has a density $\rho$, mass M and co-efficient of linear expansion $\alpha$. If heat capacities are related\par
C$_\mathrm p$ $-$ C$_\mathrm v$ = x\par
Find the value of x.
\begin{qoptions}
\opt{A}{$\dfrac23\propto\dfrac{\mathrm{MP}}{\rho}$}
\opt{B}{$\dfrac32\propto\dfrac{\mathrm{MP}}{\rho}$}
\opt{C}{$\dfrac12\propto\dfrac{\mathrm{MP}}{\rho}$}
\opt{D}{$3\propto\dfrac{\mathrm{MP}}{\rho}$}
\end{qoptions}
\qn{5} A potato gun fires a potato of length $\ell$ horizontally in a cylinder of cross-sectional area A open at one end. When the gun is fired, the potato slug is at rest, the volume between the closed end of the cylinder and the potato is V$_0$, and the pressure of the gas in this volume is P$_0$. The atmospheric pressure is P$_\mathrm{atm}$, where P$_0$ $>$ P$_\mathrm{atm}$. The gas in the cylinder is diatomic. The potato moves along the cylinder quickly enough that no heat is transferred to the gas. Friction between the potato and the barrel is negligible and no gas leaks around the potato. The parameters P$_0$, P$_\mathrm{atm}$, V$_0$, and A are fixed, but the overall length L of the barrel may be varied.(Assume $\ell<<$ L)\par
\qfig{p1_phy_s0_q05_1}
\begin{qoptions}
\opt{A}{The maximum kinetic energy E$_\mathrm{max}$ with which the potato can exit the barrel is $\left[\dfrac52\mathrm P_0+\mathrm P_\mathrm{atm}-\dfrac52(\mathrm P_\mathrm{atm})^{2/7}\mathrm P_0^{5/7}\right]\mathrm V_0$}
\opt{B}{The maximum kinetic energy E$_\mathrm{max}$ with which the potato can exit the barrel is $\left[\dfrac52\mathrm P_0+\mathrm P_\mathrm{atm}-\dfrac72\mathrm P_\mathrm{atm}^{2/7}\mathrm P_0^{5/7}\right]\mathrm V_0$}
\opt{C}{The length L in this case is $\dfrac{\mathrm V_0}{\mathrm A}\left(\dfrac{\mathrm P_0}{\mathrm P_\mathrm{atm}}\right)^{5/7}$}
\opt{D}{The length L in this case is $\dfrac{\mathrm V_0}{2\mathrm A}\left(\dfrac{\mathrm P_0}{\mathrm P_\mathrm{atm}}\right)^{5/7}$}
\end{qoptions}
\qn{6} Two cylinders having frictionless pistons A and B are connected by small tube through which a light inextensible string (which is connected to centres of both the pistons) passes. Both the pistons are having same thickness and made up of same material. Radius of piston A is twice the radius of piston B. P is the pressure of gas enclosed between the pistons and P$_0$ be the atmospheric pressure.\par
String is in horizontal plane\par
\qfig{p1_phy_s0_q06_1}
\begin{qoptions}
\opt{A}{If tension in the string is non zero both the pistons must accelerate}
\opt{B}{Both the pistons will be in equilibrium only when P = P$_0$}
\opt{C}{If string is in vertical plane such that the piston A is on the top, the equilibrium of pistons is not possible}
\opt{D}{If string is in vertical plane such that the piston B is on the top, the equilibrium of pistons is not possible}
\end{qoptions}
\qn{7} A certain mass of saturated water vapour is contained in a cylindrical vessel under a weightless piston at one atmospheric pressure. The piston is slowly lowered till 0.7 g of vapour gets condensed. The volume under the piston is now only (1/5)$^\mathrm{th}$ of the original volume. (Neglect the volume of liquid formed.) Then,
\begin{qoptions}
\opt{A}{original mass of the vapour was 0.89 gm}
\opt{B}{original mass of the vapour was 1.14 gm}
\opt{C}{temperature of vapour is 100$^\circ$C}
\opt{D}{temperature of vapour is 114$^\circ$C}
\end{qoptions}
\qn{8} Figure shows a double chambered vessel with thermally insulated walls and partitions. On each side there are n moles of an ideal monoatomic gas. Initially the pressure, volume and temperature in each side is P$_0$, V$_0$, T$_0$. The heater in first chamber supplies heat very slowly till the gas in the first chamber expands such that the pressure, volume and temperature of the gas on the left side is P$_1$, V$_1$, T$_1$ respectively and on right chamber is P$_2$ = $\dfrac{27\mathrm P_0}8$, V$_2$ and T$_2$ respectively :-\par
\qfig{p1_phy_s0_q08_1}
\begin{qoptions}
\opt{A}{Volume of first chamber is $\left[2-\left(\dfrac8{27}\right)^{3/5}\right]\mathrm V_0$}
\opt{B}{Temperature in second chamber is $\left(\dfrac{27}8\right)^{2/5}\mathrm T_0$}
\opt{C}{Work done on the gas in second chamber in terms of molar heat capacity at constant volume and T$_0$ is $\mathrm{nC}_\mathrm V\mathrm T_0\left[\left(\dfrac{27}8\right)^{2/5}-1\right]$}
\opt{D}{Work done on the gas in first chamber in terms of molar heat capacity at constant volume and T$_0$ is $\mathrm{nC}_\mathrm V\mathrm T_0\left[\left(\dfrac{27}8\right)^{2/5}-1\right]$}
\end{qoptions}
\qn{9} RMS velocity of molecules of di-atomic gas to be increased to 1.5 times. Then
\begin{qoptions}
\opt{A}{The ratio fo initial volume of final volume, if it is done adiabatically is (1.5)$^5$}
\opt{B}{The ratio of initial volume to final volume, if it is done adiabatically is (2)$^5$}
\opt{C}{The ratio of initial volume to final volume, if it is done isobarically is $\dfrac49$}
\opt{D}{The ratio of work done by gas during adiabatic to isobaric processes is $-$2.5}
\end{qoptions}
\qn{10} A thermally insulated chamber of volume 2V$_0$ is divided by a frictionless piston of area S into two equal parts A and B. Part A has an ideal gas at pressure P$_1$ and temperature T$_0$ and part B is vacuum. A massless spring of force constant K is connected with the piston and the wall of the container as shown. Initially the spring is unstretched the gas inside chamber A is allowed to expand. Let in equilibrium the spring be compressed by x$_0$. Thus :\par
\qfig{p1_phy_s0_q10_1}
\begin{qoptions}
\opt{A}{final pressure of the gas is $\dfrac{\mathrm{kx}_0}{\mathrm S}$}
\opt{B}{work done by the gas is $\dfrac12\mathrm{kx}_0^2$}
\opt{C}{change in internal energy of the gas is $\dfrac12\mathrm{kx}_0^2$}
\opt{D}{temperature of the gas is decreased}
\end{qoptions}
\qn{11} An ideal gas is taken round a cyclic thermodynamic process ABCA as shown in the given figure. If the internal energy of the gas at point A is assumed zero, while at B it is 50 J. The heat absorb by the gas in the process BC is 90 J, then\par
\qfig{p1_phy_s0_q11_1}
\begin{qoptions}
\opt{A}{The internal energy of the gas at point C is 140 J.}
\opt{B}{The heat rejected or absorbed by the gas in the process CA is 50 J}
\opt{C}{The heat absorbed by the gas in the process AB is 70 J}
\opt{D}{The net work done by the gas in the complete cycle ABCA is 20 J}
\end{qoptions}
\qn{12} One mole of an ideal gas has internal energy U which varies with volume V as U = U$_0$ + 2PV, where P = Pressure of gas, U$_0$ = constant. The gas undergoes a quasistatic cyclic process as shown in the figure (U-V diagram)\par
\qfig{p1_phy_s0_q12_1}
\begin{qoptions}
\opt{A}{Work done by the gas in the process A to B is $\left(\dfrac{\mathrm U_2-\mathrm U_1}2\right)\ln3$}
\opt{B}{Work done by the gas in the process A to B is $\left(\dfrac{\mathrm U_2-\mathrm U_0}2\right)\ln3$}
\opt{C}{The gas is mixer of mono and diatomic}
\opt{D}{The gas is mixer of di- and polyatomic}
\end{qoptions}
\qn{13} A sample of ideal gas is taken through the cyclic process shown in the figure. The temperature of the gas at state A is T$_\mathrm A$ =200 \ K. At states B and C the temperature of the gas is the same. The greatest temperature of the gas in kelvin during the cyclic process is T$_\mathrm{max}$. Find T$_\mathrm{max}$/800.\par
\qfig{p1_phy_s0_q13_1}
\qn{14} Compute the amount of work done by a gas which is being isobarically heated from 27$^\circ$C to 127$^\circ$C if it is in a vessel closed by a movable piston with a cross-sectional area of 20 cm$^2$ and weighing 5 kg. Consider two cases: (1) the vessel is arranged horizontally (Figure-a), and (2) the vessel is arranged vertically (Figure-b). The initial volume of the gas is 6 litres. Assume standard atmospheric pressure = 100 kPa. Neglect friction.\par
Fill (w$_1$ + w$_2$)/50 \ (in J) in OMR.\par
\qfig{p1_phy_s0_q14_2}
\qn{15} One mole of an ideal monoatomic gas is kept in a rigid vessel whose outer surface is maintained at a constant temperature 100$^\circ$C. The thickness and surface area of the vessel are 5mm and 400 cm$^2$ respectively. Initial temperature of the gas in the vessel is 16$^\circ$C. The conductivity of the material of the walls is 0.01 Wm$^{-1}$ K$^{-1}$. A sound signal move with a velocity 320 + n in the gas contained in the vessel after 250 seconds if initially its speed in the gas is 290 m/s. Find the value of n in approximate integer.\par
(Given e1.6 = 4.94, gas constant R = $\dfrac{25}3$ J mol$^{-1}$ K$^{-1}$. (Neglect the heat capacity of the vessel).
\qn{16} An ideal mono-atomic gas is trapped in a vertical fixed cylinder under an airtight piston of mass m and area S. Initially, when the piston is held at height h$_0$ above the bottom of the cylinder, the pressure of the gas is P$_0$. Outside the cylinder is vacuum and heat loss to the surroundings is negligible. When the piston is released, it rises \ and acquires its maximum speed after moving a distance h. If maximum speed of piston is\par
$\mathrm V=\sqrt{\dfrac{x\mathrm P_0\mathrm{Sh}_0}{\mathrm m}-\mathrm g(5\mathrm h+y\mathrm h_0)}$ . Find (x + $y$). Acceleration of free fall is g.\par
\qfig{p1_phy_s0_q16_1}
\qn{17} The value of $\gamma=\dfrac{\mathrm C_\mathrm p}{\mathrm C_\mathrm v}$ is $\dfrac43$ for an adiabatic process of any ideal gas for which internal energy U = K + nPV. The value of n (K is a constant) is :
\qn{18} Between two isotherm we have a cycle as shown. The work done by the gas during the cycle is 15 n (in J). Find the value of n. (Take T$_1$ = 127$^\circ$C, T$_2$ = 16$^\circ$C, n = 1 mole]\par
\qfig{p1_phy_s0_q18_1}
\raggedcolumns
\papersec{1}
\partsec{2}{CHEMISTRY}
\qn{19} Which of the following statement is false?
\begin{qoptions}
\opt{A}{Strontium decomposes water more readily than beryllium}
\opt{B}{Barium carbonate melts at a higher temperature than calcium carbonate}
\opt{C}{Stable superoxides of all alkali metals exist and are paramagnetic}
\opt{D}{Sodium ions are discharged in preference to hydrogen ions at a mercury cathode.}
\end{qoptions}
\qn{20} \textbf{Consider the following statements,}\par
\textbf{S$_1$} : Ozone reacts with dry KOH forming a deep red coloured compound and liberating oxygen gas at low temperature.\par
\textbf{S$_2$} : KO$_2$ is better oxidant then K$_2$O$_2$ both in aqueous and acidic medium.\par
\textbf{S$_3$} : KOH is better absorber of CO$_2$ than NaOH.\par
\textbf{S$_4$} : Quick lime on heating with carbon in an electric furnace forms CaCO$_3$.\par
and arrange in the order of true/false.
\begin{qoptions}
\opt{A}{T T T F}
\opt{B}{T T F T}
\opt{C}{F F T F}
\opt{D}{F T T F}
\end{qoptions}
\qn{21} Silicate having one monovalent corner oxygen atom in each tetrahedron unit is :
\begin{qoptions}
\opt{A}{sheet silicate}
\opt{B}{cyclic silicate}
\opt{C}{single chain silicate}
\opt{D}{double chain silicate}
\end{qoptions}
\qn{22} Heating a mixture of ammonium chloride and sodium tetrahydridoborate gives one liquid product (X), along with other products, under ambient conditions. Compound (X) is :
\begin{qoptions}
\opt{A}{NH$_4$[BH$_4$]}
\opt{B}{[(NH$_3$)$_2$BH$_2$][BH$_4$]}
\opt{C}{N$_3$B$_3$H$_6$}
\opt{D}{N$_3$B$_3$H$_{12}$}
\end{qoptions}
\qn{23} A compound `X' on heating gives a colourless gas which is not supporter of combustion. The residue is dissolved in water and excess of SO$_2$ is passed through it. First we get while ppt which dissolves in excess of SO$_2$. The compound `X' cannot be
\begin{qoptions}
\opt{A}{CaCO$_3$}
\opt{B}{NaNO$_3$}
\opt{C}{CaSO$_4$.2H$_2$O}
\opt{D}{K$_2$CO$_3$}
\end{qoptions}
\qn{24} In the structure of borax :
\begin{qoptions}
\opt{A}{5 B--O--B bonds are present}
\opt{B}{The number of $\sigma$ bonds is 34}
\opt{C}{19 atoms are sp$^3$ and another 2 atoms are sp$^2$ hybridized}
\opt{D}{The number of lone pair's are 34}
\end{qoptions}
\qn{25} Which of the following statements is/are \textbf{CORRECT} w.r.t Boron for its compounds.
\begin{qoptions}
\opt{A}{Boric acid with NaOH does not give a sharp end point but when taken along with sugars it gives a sharp end point with NaOH.}
\opt{B}{Aqueous solution of Borax shows buffer properties.}
\opt{C}{B$_2$H$_6$ on reaction with excess NH$_3$ at low temperatures forms an ionic compound which on heating gives an aromatic compound which is less reactive than benzene.}
\opt{D}{On hydrolysis of BF$_3$, boric acid along with flouroborates are obtained.}
\end{qoptions}
\qn{26} When a solution of alkaline earth metals in liquid ammonia is evaporated :
\begin{qoptions}
\opt{A}{They form hexa ammoniates [M(NH$_3$)$_6$]}
\opt{B}{The ammonia decomposes at high temperature}
\opt{C}{The tendency to form ammoniates decreases with increase in the size}
\opt{D}{The ammoniates are good conductors of electricity.}
\end{qoptions}
\qn{27} Ca$_2$B$_6$O$_{11}$ + Na$_2$CO$_3$ $\xrightarrow{\Delta}$[X] + CaCO$_3$ + NaBO$_2$ \ (unbalanced equation)
\begin{qoptions}
\opt{A}{structure of anion of crystalline [X] has one boron atom sp$^3$ hybridised and other three boron atoms sp$^2$ hybridised.}
\opt{B}{X with NaOH (aq.) gives a compound which on reaction with hydrogen peroxide in alkaline medium yields a compound used as brighter in soaps.}
\opt{C}{hydrolysis of [X] with HCl or H$_2$SO$_4$ yields a compound which on reaction with HF gives fluoroboric acid}
\opt{D}{[X] on heating with chromium salts in oxidising flame gives green coloured bead in cold.}
\end{qoptions}
\qn{28} Consider the following reactions of Magnesium, and select incorrect statements in context to three products X, Y, And Z \ ``only''\par
3Mg + N$_2$ $\xrightarrow{\text{heat}}$ ``X''\par
2Mg + CO$_2$ $\to$ ``Y'' + \ C\par
2Mg + SO$_2$ $\to$ ``Y'' + S\par
Mg + 2HNO$_3$ $\to$ ``Z'' + \ H$_2$
\begin{qoptions}
\opt{A}{Maximum \% \ magnesium is in compound Z}
\opt{B}{Maximum \% magnesium is in compound Y}
\opt{C}{X can also be produced during reaction of Mg with NH$_3$ at high temperature.}
\opt{D}{Aq. solution of \ Y and Z both are ``basic''}
\end{qoptions}
\qn{29} Select correct ``reaction and observation'' combination(s).\par
\begin{tabular}{@{}l p{3.4cm} p{3.4cm}@{}}
 & \underline{\textbf{Reactions}} & \underline{\textbf{Observations}}\\
(i) & NaCl + NH$_4$HCO$_3$ $\longrightarrow$ Product(s) & Polyatomic anion in product species has sp$^3$ hybrid central atom\\
(ii) & Na$_2$CO$_3$ + Ca(OH)$_2$ $\longrightarrow$ Product(s) & One of the products is insoluble precipitate\\
(iii) & KOH + O$_3$ $\longrightarrow$ Product(s) & Products include - a neutral oxide, two paramagnetic species\\
\end{tabular}\par
\begin{qoptions}
\opt{A}{(i) , (ii) only}
\opt{B}{(ii) , (iii) only}
\opt{C}{(ii) only}
\opt{D}{(i) , (ii) \& (iii)}
\end{qoptions}
\qn{30} Silicones are a group of organo silicon polymer, select correct statement(s) about them.
\begin{qoptions}
\opt{A}{R$_3$SiCl produce very complex cross linked polymeric silicons on hydrolysis.}
\opt{B}{Their water rapellency arises because a silicone chain is surrounded by organic side groups.}
\opt{C}{Their strength and interness are related to their stable silica like skeleton of $\begin{array}{*{11}{c}} & | &  &  &  & | &  &  &  & | & \\{-} & \mathrm{Si} & {-} & \mathrm O & {-} & \mathrm{Si} & {-} & \mathrm O & {-} & \mathrm{Si} & {-}\\ & | &  &  &  & | &  &  &  & | & \end{array}$}
\opt{D}{Hydrolysis product of CH$_3$SiCl$_3$ will block the end of the straight chain produce by hydrolysis of (CH$_3$)$_2$SiCl$_2$}
\end{qoptions}
\qn{31} Identify the no. of correct statements in context to given metals ``only''\par
Li, Na, K, Mg, Be, Sr and Ba\par
(i) Three elements on heating with nitrogen form stable nitrides.\par
(ii) Only two alkaline earth metals reacts with cold water easily/rapidly.\par
(iii) Diagonal relationship pair is presnt in the given list of metals.\par
(iv) Densest alkali metal is in above list.\par
(v) Two of them create amphoteric oxide.\par
(vi) Highest Boiling point is of Mg.\par
(vii) All impart color to flame in flame coloration test.
\qn{32} Select the no. of correct statements\par
(i) \ Fusion mixture is (Na$_2$CO$_3$ + K$_2$CO$_3$) and black ash contains ( Na$_2$CO$_3$ and CaS).\par
(ii) Sodium Bicarbonate is obtained as an intermediate product in Solvay ammonia process.\par
(iii) On heating, Sodium Bicarbonate at around 500K, we get basic oxide, acidic oxide and neutral oxide.\par
(iv) Mixture of sodium bicarbonate, starch, and potassium hydrogen tartarate is known as \textit{baking powder}.\par
(v) Sea water contains KCl as well as NaCl.\par
(vi) Hot concentrated ``sylvine'' on cooling gives out the crystals of KCl first, NaCl being more soluble remains in the solution.\par
(vii) KOH can be prepared by the electrolysis of an aqueous solution of potassium chloride.\par
(viii) Mg(OH)$_2$ is used as an antacid.\par
(ix) A mixture of MgO and MgCl$_2$ having composition MgCl$_2$. 5MgO. xH$_2$O is called Sorel's cement
\qn{33} Find the sum of number of pentagon (x), number of hexagon (\textit{y}) \& number of Carbon--Carbon linkage (Z) in fullerene(C$_{60}$)
\qn{34} A comparative study of following properties between benzene and borazine are taken. For how many properties data of benzene is more than borazine.\par
Melting point, Boiling point, Dipole moment, Bond polarity, $|$Resonance energy$|$, Molecular weight.
\qn{35} Find the number of \textbf{CORRECT} statements about B$_2$H$_6$\par
(i) It is lewis acidic gas\par
(ii) It gives symmetrical bond cleavage with CO\par
(iii) It gives H$_2$ gas with C$_2$H$_5$OH\par
(iv) It gives borazene with NH$_3$ on mild heating\par
(v) It hydrolysed easily with H$_2$O and gives boric acid\par
(vi) All the adjacent bond angles are identical to each other
\qn{36} Number of B--B bonds = x\par
Number of B--H--B bonds = \textit{y}\par
Find the value of (x + \textit{y}) in B$_4$H$_{10}$.
\raggedcolumns
\papersec{1}
\partsec{3}{MATHEMATICS}
\qn{37} Number of ways in which 4 married couples can be seated round a table such that men and women are alternate and not all women adjacent to their husband, is
\begin{qoptions}
\opt{A}{144}
\opt{B}{132}
\opt{C}{120}
\opt{D}{36}
\end{qoptions}
\qn{38} If x, y, z, w $\in$ N, such that each of them gives remainder 3 when divided by 5, then number of solutions of the equation \ x + y + z + w = 62, is equal to
\begin{qoptions}
\opt{A}{$^{13}$C$_3$}
\opt{B}{$^9$C$_3$}
\opt{C}{$^{12}$C$_3$}
\opt{D}{$^{12}$C$_2$}
\end{qoptions}
\qn{39} The number of ways in which 5 dance couples (consisting of 1 man and 1 woman in each dance couple) can be formed from 5 married couples (including Naresh and Rani) such that no husband dances with his own wife but Naresh wants to dance with Rani (not his wife)
\begin{qoptions}
\opt{A}{44}
\opt{B}{11}
\opt{C}{65}
\opt{D}{205}
\end{qoptions}
\qn{40} Four digit numbers are formed using the digits 1, 2, 3, 4 (repetition is allowed). The number of such four digit numbers divisible by 11 is :
\begin{qoptions}
\opt{A}{22}
\opt{B}{28}
\opt{C}{44}
\opt{D}{none of these}
\end{qoptions}
\qn{41} A shopkeeper places before you 41 different toys out which 20 toys are to be purchased. Suppose m \ = \ number of ways in which 20 toys can be purchased without any restriction and \ n = number of ways in which a particular toy is to be always included in each selection of 20 toys, \ then (m $-$ n) can be expressed as :
\begin{qoptions}
\opt{A}{$\dfrac{2^{10}}{20!}(1\cdot3\cdot5\dots\dots\dots\cdot39)$}
\opt{B}{$\dfrac{2^{20}(1\cdot3\cdot5\dots\dots\dots\cdot19)}{10!}$}
\opt{C}{$\displaystyle\prod_{\mathrm r=0}^{19}\left(\dfrac{4\mathrm r+2}{20-\mathrm r}\right)$}
\opt{D}{$\binom{21}{1}\binom{22}{2}\binom{23}{3}\dots\dots\binom{40}{20}$}
\end{qoptions}
\qn{42} Consider the word \ Z = F R E E W H E E L. Which of the following statement(s) is/are correct
\begin{qoptions}
\opt{A}{Number of other ways in which the letters of the word Z can be arranged is $^9$P$_5$.}
\opt{B}{Number of ways in which the letters of the word Z can be an arranged in a circle is \ $^8$P$_4$ distisguishing between clockwise and anticlockwise.}
\opt{C}{Number of ways in which the letters of the word Z can be arranged if vowels and consonants both are in alphabetical order is $\dfrac12\,{}^{10}\mathrm C_5$}
\opt{D}{If as many more words are formed as possible using the letters of the word Z the number of words which contain the word FEEL is $^6$P$_4$.}
\end{qoptions}
\qn{43} Number of ways in which 6 different toys can be distributed equally among \ 3 \ children is also equal to
\begin{qoptions}
\opt{A}{number of ways in which all the letters of the word \ \ KUMKUM \ \ can be arranged in a definite order.}
\opt{B}{number of different signals which can be obtained from 10 different flags taken two at a time keeping one above the other.}
\opt{C}{number of ways in which five different marbles can be placed in three different boxes so that no box remains empty.}
\opt{D}{number of ways in which all the letters of the word CALCULUS can be arranged without changing the order and position of \ non alike letters.}
\end{qoptions}
\qn{44} A multiple choice question in an examination consists of \ 5 alternatives in which one or more than one may be correct. Number of different possible answers that can be given by an examinee is also equal to
\begin{qoptions}
\opt{A}{The number of all possible forecasts of 5 football matches in which atleast one forecast is incorrect if each forecast is either a win or loss.}
\opt{B}{The value of (p + q) where p and q satisfy the equations C (p, p $-$ 2) = 10 \ and \ C (q, 15) = C (q, 11).}
\opt{C}{Number of ways in which a normal dice tossed twice can fall in a definite order given that both throws are different.}
\opt{D}{Number of different signals that can be made with 5 different flags by raising any number of flags at a time.}
\end{qoptions}
\qn{45} Let P$_\mathrm n$ denotes the number of permutations of n distinct things taken all at a time and\par
$\mathrm x_\mathrm n={}^{\mathrm n+5}\mathrm C_4-\left(\dfrac{143}{96}\right)\left(\dfrac{\mathrm P_{\mathrm n+5}}{\mathrm P_{\mathrm n+3}}\right)$\par
(where n $\in$ N). The possible value of \ n \ for which \ x$_\mathrm n$ is negative, can be
\begin{qoptions}
\opt{A}{1}
\opt{B}{2}
\opt{C}{3}
\opt{D}{4}
\end{qoptions}
\qn{46} A man has 3 friends. Number of ways in which he can invite one friend every day for dinner on 6 successive nights so that he invites every friend atleast once is N. Which of the following hold(s) good?
\begin{qoptions}
\opt{A}{Number of divisors of \ N which are divisible by 5 is 12.}
\opt{B}{Sum of all the divisors of N is 1680.}
\opt{C}{Number of ways in which \ N can be resolved as product of two divisors which are relatively prime, is 4.}
\opt{D}{Number of divisors of the form (4n + 2), n $\in$ N is 6.}
\end{qoptions}
\qn{47} All the 7-digit Numbers containing each of the digits 1, 2, 3, 4, 5, 6, 7 exactly once, and not divisible by 5, are arranged in increasing order. Then which of the following is/are CORRECT ?
\begin{qoptions}
\opt{A}{1993$^\mathrm{rd}$ number in the list is 4312567}
\opt{B}{1996$^\mathrm{th}$ number in the list is 4312756}
\opt{C}{2000$^\mathrm{th}$ number in the list is 4315672}
\opt{D}{1999$^\mathrm{th}$ number in the list is \ 43152763}
\end{qoptions}
\qn{48} If a natural number, \ N = 2$^3$ $\cdot$ 3$^2$ $\cdot$ 5$^1$ $\cdot$ 7$^1$ \ then identify which of the following statement(s) is(are) correct?
\begin{qoptions}
\opt{A}{The product of all divisors of N is (N)$^{24}$}
\opt{B}{The product of all even divisors of \ N is \ (2N)$^{18}$}
\opt{C}{The product of all odd divisors of \ N is \ $\left(\dfrac{\mathrm N}2\right)^6$}
\opt{D}{The product of all composite divisors of N is \ $\left(\dfrac{\mathrm N^{24}}{210}\right)$}
\end{qoptions}
\qn{49} There are 5 white, 4 yellow, 3 green, 2 blue and 1 red ball. The balls are all identical except for colour. These are to be arranged in a line in 5 places. Find the number of distinct arrangements.
\qn{50} In a league of 8 teams, each team played every other team 10 times. The number of wins of the 8 teams formed an arithmetic sequence. Find the least possible number of games won by the champion.
\qn{51} Find the number of \ 4-digit numbers with distinct digits chosen from the set \{0, 1, 2, 3, 4, 5\} in which no two adjacent digits are even.
\qn{52} If the number of ways in which \ 14 identical toys can be distributed among three boys so that each one gets at least one toy and no two boys get equal number of toys is m, then find the value of $\left(\dfrac m{20}\right)$.
\qn{53} W$_1$, W$_2$ \ and \ W$_3$ \ are three ticket windows in a cinema hall. 12 persons \ (A$_1$, A$_2$, \dots\dots, A$_{12}$) \ are standing for the tickets in a queue of \ 4, 3 \ and 5 number of persons infront of \ W$_1$, W$_2$ \ and W$_3$ \ respectively. A$_1$ \ and A$_2$ \ want to get their tickets from \ W$_1$ \ and A$_3$, A$_4$ \ and \ A$_5$ want to get their tickets from \ W$_3$. If the number of ways in which all are getting their tickets such that \ A$_1$, A$_2$, A$_3$, A$_4$ \ and \ A$_5$ \ are getting their tickets as early as possible is \ k(7!), then find the value of $\left[\dfrac{\mathrm k}3\right]$.\par
(Assuming each person is getting exactly one ticket at a time ).\par
[\textbf{Note:} [y] denotes greatest integer less than or equal to \textit{y}.]
\qn{54} Let f : A $\to$ B \ be any \ function \ where A \ is a set containing the positive integral solution of the inequality \ $\mathrm{cosec}^{-1}(\mathrm{cosec}\,2)>x^2-3x$ and B is the set of all divisors of the natural number 2010. If \ f(i)$\le$ f(j) \ $\forall$ i $<$ j, then find the total number of mappings \ from A \ to B.
\end{multicols}
\end{document}
